
\begin{table}[htbp]
   \caption{\label{tab:category_decomp} Category-level robustness: extensive vs.\ intensive decomposition of the rise in quantity sold}
   \centering
   \begin{tabular}{lccc}
      \tabularnewline \midrule \midrule
      Model:          & (1)           & (2)           & (3)\\  
      \midrule
      \emph{Variables}\\
      SOE$_{st}$      & 0.110$^{***}$ & 0.080$^{***}$ & 0.029$^{***}$\\   
                      & (0.010)       & (0.009)       & (0.001)\\   
      Post-SOE$_{st}$ & 0.017$^{*}$   & 0.003         & 0.014$^{***}$\\   
                      & (0.008)       & (0.010)       & (0.002)\\   
      \midrule
      \emph{Fixed-effects}\\
      category        & Yes           & Yes           & Yes\\  
      store\_id       & Yes           & Yes           & Yes\\  
      \midrule
      \emph{Fit statistics}\\
      Observations    & 629,486       & 629,486       & 629,486\\  
      R$^2$           & 0.79219       & 0.84983       & 0.91996\\  
      \midrule \midrule
      \multicolumn{4}{l}{\emph{Clustered (sst) standard-errors in parentheses}}\\
      \multicolumn{4}{l}{\emph{Signif. Codes: ***: 0.01, **: 0.05, *: 0.1}}\\
   \end{tabular}
   
   \par \raggedright 
   Aggregated over ALL UPCs in each focal category (store-category-week grain).\\
   Column (1): log category volume. Column (2): log purchase occasions (extensive margin).\\
   Column (3): log volume per occasion (intensive margin). $\beta_Q = \beta_N + \beta_{Q/N}$.\\
   The occasion count sums per-variety transaction counts, so a basket buying two\\
   varieties of a category counts twice; this modestly inflates the extensive margin.\\
   FEs: category and store. Standard errors clustered at the state level.
\end{table}


